\documentclass[10pt,a4paper]{amsart}
\usepackage[margin=19mm]{geometry}
\usepackage{amsmath,amssymb,mathtools}
\usepackage[scr=rsfs,cal=euler]{mathalfa}
\usepackage{xfrac}
\usepackage[unicode,hidelinks,bookmarks=false]{hyperref}
\newcommand{\ie}{i.e.\ }
\newcommand{\cf}{cf.\ }
\newcommand{\resp}{respectively\ }
\newcommand{\Subsection}{\subsection}
\providecommand{\nobreakdash}{\nobreak\hbox{-}}
\setlength{\emergencystretch}{2em}
\multlinegap0pt
\usepackage{amssymb}
\usepackage{xcolor}
\usepackage[shortlabels]{enumitem}
\newtheorem{theorem}{Theorem}
\usepackage{cleveref}
\crefname{theorem}{Theorem}{Theorems}
\def\COMMENT{}
\newcommand{\F}{\mathcal{F}}
\def\ev{\kern-1pt\lowbkwd e{0.5}2\kern-1pt}
\def\fv{\lowbkwd f01}
\def\fwd #1#2{{\lowfwd{#1}{#2}{15}}}
\renewcommand{\geq}{\geqslant}
\def\lowbkwd #1#2#3{{\mathop{\kern0pt #1}\limits^{\kern#2pt\raise.#3ex
\vbox to 0pt{\hbox{$\scriptscriptstyle\leftarrow$}\vss}}}}
\def\lowfwd #1#2#3{{\mathop{\kern0pt #1}\limits^{\kern#2pt\raise.#3ex
\vbox to 0pt{\hbox{$\scriptscriptstyle\rightarrow$}\vss}}}}
\def\vE{{\hskip-1pt{\fwd E3}\hskip-1pt}} %{{\vec S}} %
\def\vF{{\hskip-1pt{\fwd F3}\hskip-1pt}} %{{\vec S}} %
\def\vS{{\hskip-1pt{\fwd S3}\hskip-1pt}} %{{\vec S}} %
\def\ve{\kern-1.5pt\lowfwd e{1.5}2\kern-1pt}
\def\vf{\kern-2pt\lowfwd f{2.5}2\kern-1pt}
\def\vs{\hskip-1pt\lowfwd s{1.5}1}
\title{Tangle structure trees~II:\\ trees of tangles and tangle-tree duality\hbox to 3mm{\hfil}}
\author{Hanno von Bergen, Reinhard Diestel}
\date{Source: arXiv:2603.17756v1 (CC BY 4.0)}
\begin{document}
\maketitle

\noindent\emph{Source and licence.} Tangle structure trees~II:\\ trees of tangles and tangle-tree duality\hbox to 3mm{\hfil}. Version arXiv:2603.17756v1 (CC BY 4.0), distributed by arXiv under the Creative Commons Attribution 4.0 International licence, http://creativecommons.org/licenses/by/4.0/, as recorded in the arXiv licence metadata for this exact version and shown on the abstract page https://arxiv.org/abs/2603.17756v1. Full licence notice: \texttt{licenses/arxiv-2603.17756-CC-BY-4.0.txt}.\par\smallskip

\noindent\textit{Editorial scope.} Counted unit: the article's theorem thm:main (source0.tex lines 541--552). The article's complete original proof of it is delivered with it (source0.tex lines 570--613, one \texttt{proof} environment of 44 lines, which the article places away from the statement). No mathematics is written, repaired or re-derived: the statement and the proof are the article's own lines, in the article's own notation, and no retained line is substituted or re-flowed.  No further source-local context is needed: every cross-reference of every retained line resolves inside the two delivered spans.  Nothing was omitted from any retained span, so the package carries no omission record.  No retained line contains a citation, so the wrapper carries no bibliography and no bibliography entry.  No retained line contains a figure environment, an image-inclusion command or drawing code, so no image is delivered and the package declares no asset dependency.  Editorial formatting only, in addition: an AMS \texttt{amsart} preamble that restates the article's own declarations and macros used by the retained lines, namely the environments \texttt{theorem} on separate counters, together with the standard packages those lines need.  The retained environments print in this document as Theorem 1 in order of appearance, whereas the article prints the same items with the numbering of its own full document.  The complete native source of the article as distributed in the arXiv e-print for this exact version is bundled unchanged as \texttt{sources/196645/source0.tex}.
\begin{theorem}\label{thm:main}
Let $\vS$ be a separation system without trivial elements, ${\F\subseteq 2^{\vS}}\!\!$ a~set of stars, and $(T,r,\beta)$ an irreducible $\F$-tree of~$\vS$. Then there exists an $S$-tree $(T', \alpha)$ over $\F$ and a map $\gamma\colon V(T') \cup \vE(T') \to V(T) \cup E(T)$ such that%
   \COMMENT{}
\begin{enumerate}\itemsep2pt\vskip4pt
    \item\label{item:main1} $\gamma$ is a bijection from the nodes of $T'$ to the leaves of $T$;
    \item\label{item:main2} $\gamma$ is a bijection from the oriented edges of $T'$ to the edges of~$T$;
    \item\label{item:main3} for every edge~$e$ of~$T'\!$ there is a node~$v_e$ of~$T\!$ such that $\{\gamma(\ve), \gamma(\ev)\} = E_{v_e}$, and these~$v_e$ are distinct for different~$e$;
    \item\label{item:main4} for every oriented edge $\ve$ of $T'$, the map $\gamma$ sends $t(\ve)$ to a leaf $\ell > v_e$ of $T$ such that $\gamma(\ve)$ is the first edge of the path $v_eT\ell$;
    \item\label{item:main5} $\alpha = \beta \circ \gamma$.
\end{enumerate}\vskip2pt
In particular, $\beta(E(T)) = \alpha(\vE(T'))$. 
\end{theorem}

\begin{proof}[Proof of \cref{thm:main}]
We start our construction of~$T'$ by taking as~$V(T')$ any set we can map to the leaves of~$T$ by a bijection~$\gamma$. This ensures~\ref{item:main1}.

We shall give~$T'$ one edge~$e$ for every non-leaf node $v =: v_e$ of~$T$. To choose the ends of~$e$ in~$T'$, consider the children $w_1$ and~$w_2$ of~$v$ in~$T$.%
   \footnote{These are distinct: $s_v$~cannot be degenerate, because $vw_1$ is necessary for a forbidden%
   \COMMENT{}
   leaf $\ell\ge w_1$ of~$T$, so $\beta(vw_1)$ lies in a star $\sigma\in\F$. Stars do not contain degenerate separations.}%
   \COMMENT{}
   For each $i\in\{1,2\}$, since $T$ is irreducible, the edge~$vw_i$ is necessary for some leaf~$\ell_i\ge w_i$ of~$T$.%
   \footnote{Our proof will imply that these~$\ell_i$ are unique, but for the definition it suffices to pick any.}%
   \COMMENT{}
   Let $e$ join the nodes $\gamma^{-1}(\ell_1)$ and~$\gamma^{-1}(\ell_2)$ of~$T'$, and let $\gamma(\ve) := vw_i$ for the orientation~$\ve$ of~$e$ for which $\gamma(t(\ve)) = \ell_i$. This satisfies \ref{item:main2}--\ref{item:main4}, if we think of~$T'$ as a multigraph for now, i.e., allow multiple edges.

However, since $v_e$ is the unique infimum of~$\ell_1$ and~$\ell_2$ in~$T$, the pair $\gamma^{-1}(\{\ell_1,\ell_2\})$ of nodes in~$T'$ is joined only by the edge~$e$: the graph~$T'$ has no parallel edges and no loops.
Finally, now that $\gamma$ is fixed, let $\alpha:=\beta\circ\gamma$ as in~\ref{item:main5}.

It remains to show that $T'$ is a tree, and that $(T',\alpha)$ is over~$\F$: as $\alpha(\ve)^* = \alpha(\ev)$ by definition of $\gamma$ and~$\alpha$, it will be clear that $(T',\alpha)$ is an $S$-tree once we know that $T'$ is a tree.

Let us begin by showing that $T'$ is acyclic. Consider two edges $e, f$ of $T'$ with $t(\ve) = i(\vf)$. By~\labelcref{item:main4}, the edges $\gamma(\ve)$ and $\gamma(\fv)$ of~$T$ at~$v_e$ and~$v_f$, respectively, both lie on the path~$rT\ell$ for $\ell=\gamma(t(\ve))=\gamma(i(\vf))$. As $e \neq f$, the nodes~$v_e$ and~$v_f$ of~$rT\ell$ are distinct,%
   \COMMENT{}
   by~(iii), so $v_e < v_f$ or $v_f < v_e$. As $s_u \neq s_v$ for $u < v$ in any separation tree, $\alpha(\ve)=\beta(\gamma(\ve))$ and $\alpha(\fv) = \beta(\gamma(\fv))$, compare~\labelcref{item:main5}, are thus orientations of distinct separations $s_{v_e}\!\ne s_{v_f}\!$ in~$S$.

Since we chose~$\gamma^{-1}(\ell)$ as~$t(\ve)$ and as~$t(\fv)$ when we picked the ends of the edges~$e$ and~$f$ in~$T'$, the edges $\gamma(\ve)$ and $\gamma(\fv)$ of~$T$ are necessary for~$\ell$: their $\beta$-values $\alpha(\ve)$ and $\alpha(\fv)$ each lie in some $\sigma \in \F$ that is a subset of~$\beta_{\ell}$, and every subset of~$\beta_{\ell}$ in~$\F$ contains both. Consider any such $\sigma\in\F$. Since $\sigma$ is a star, we have $\alpha(\ve) \geq \alpha(\fv)^* = \alpha(\vf)$. As $s_{v_e}\!\ne s_{v_f}$, the above inequality is strict. The $\alpha$-values along any oriented path in~$T'$ thus strictly decrease.%
   \COMMENT{}
   Hence $T'$ contains no (oriented) cycles.

To complete our proof that $T'$ is a tree, it suffices to show that it has one more node than edges. The first of these numbers equals the number of leaves of $T$; the second equals the number of non-leaves of $T$. This latter number is one less than the former, since $T$ is a binary tree. 

To show that $(T', \alpha)$ is an $S$-tree over~$\F$, consider any node~$t$ of~$T'$; we have to show that $\alpha(\vF\!_t) \in \F$. As earlier,%
   \COMMENT{}
   all the edges $\ve\in\vF\!_t$ are such that $\gamma(\ve)\in rT\ell$ for $\ell := \gamma(t)$ is necessary for~$\ell$, so $\alpha(\vF\!_t) = \beta(\gamma(\vF\!_t)) \subseteq\sigma\subseteq\beta_\ell$ for some $\sigma\in\F$. It remains to show that $\alpha(\vF\!_t)\supseteq \sigma$.

Suppose there exists an $\vs \in \sigma\setminus \alpha(\vF\!_t)$. As $\vs\in\sigma\subseteq\beta_\ell = \beta(rT\ell)$, we can find an edge $vw$ in~$rT\ell$ with $v<w$ such that $\beta(vw) = \vs$. Let $\vf := \gamma^{-1}(vw)$. Then $\alpha(\vf)=\vs$, and $\gamma(\vf)\in E(rT\ell)$ but $\gamma(\vf)\notin\gamma(\vF\!_t)\subseteq E(rT\ell)$. In particular, $v_f\ne v_e$ for every $\ve\in\vF\!_t$, and hence $f\ne e$, so neither $\vf$ or~$\fv$ lies in~$\vF\!_t$. But note that $v_f$ and all these~$v_e$ lie on the path~$rT\ell$, so they are comparable in the tree-order on~$V(T)$.

Since $T'$ is a tree, either $\vf$ or $\fv$ points towards~$t$ in~$T'$. Suppose first that $\vf$ does. Then $\vf\ge \ve$%
   \COMMENT{}
   in~$\vE(T')$ for some $\ve \in \vF\!_t$. As in our acyclicity proof earlier, this implies that $\alpha(\vf) \geq \alpha(\ve)$.%
   \COMMENT{}
   As both $\vs = \alpha(\vf)$ and~$\alpha(\ve)$ lie in the star~$\sigma$, we have $\alpha(\vf)\ge\alpha(\ve)^*$ too.%
   \COMMENT{}
   As $v_f$ and~$v_e$ are distinct nodes on~$rT\ell$, which implies that $s=s_{v_f}\ne s_{v_e} = \{\alpha(\ve),\alpha(\ve)^*\}$ since $T$ is a separation tree, $\vs = \alpha(\vf)$ is thus trivial in~$\vS$ witnessed by~$s_{v_e}$, a~contradiction to our assumptions about~$\vS$.

Suppose now that the edge~$\fv$ points towards~$t$ in~$T'$. Then $\fv \geq \ve\in\vF\!_t$ and hence $\alpha(\vf)^*\ge\alpha(\ve)$, while $s=s_{v_f}\ne s_{v_e} = \{\alpha(\ve),\alpha(\ve)^*\}$, as earlier. This makes $\vs = \alpha(\vf)$ and $\alpha(\ve)$ inconsistent elements of $\beta_{\ell}$, a~contradiction to the fact that $T$ is a consistent separation tree.
\end{proof}

\end{document}
